\chapter[Regularność De Giorgiego--Nasha--Mosera]{Teoria De Giorgiego--Nasha--Mosera: regularność przy mierzalnych współczynnikach}
\label{chap:pde17}

W rozdziałach~\ref{chap:pde15}--\ref{chap:pde16} uzyskiwaliśmy regularność
pochodnych przy odpowiedniej gładkości współczynników. Teraz przyjmujemy
\emph{jedynie mierzalność} macierzy $A$. Mimo to rozwiązanie słabe
jednostajnie eliptycznego równania skalarnego ma lokalnie reprezentanta
Hölderowskiego. Jest to szczególny przypadek twierdzenia
De Giorgiego--Nasha--Mosera. Przedstawimy metodę iteracji Mosera,
od oszacowań energetycznych aż do nierówności Harnacka i regularności.
Nie będziemy utożsamiać tej teorii z wyższą regularnością $W^{2,p}$:
przy skokowych współczynnikach pochodna rozwiązania może mieć skok.

\begin{uwaga}[Mapa założeń: zasada maksimum, Schauder i Harnack]
\label{rem:pde-v43-regularity-map}
Klasyczne zasady maksimum z rozdziału~\ref{chap:pde14}
wykorzystują regularność klasyczną albo dodatkowe argumenty
aproksymacyjne. Oszacowania Schaudera z
rozdziału~\ref{chap:pde16} wymagają regularności H\"olderowskiej
współczynników i dotyczą pochodnych drugiego rzędu.
W niniejszym rozdziale pracujemy z równaniem \emph{skalarnym}
w postaci dywergencyjnej, z mierzalnym, jednostajnie
eliptycznym $A$; regularność H\"olderowską otrzymujemy
wyłącznie dla samego rozwiązania. Nie przenosimy automatycznie
nierówności Harnacka na układy ani na równania z niekontrolowanymi
wyrazami niższego rzędu.
\end{uwaga}

Przyjmujemy oznaczenia $u_+=\max\{u,0\}$,
$\operatorname*{ess\,osc}_E u=
\operatorname*{ess\,sup}_E u-\operatorname*{ess\,inf}_E u$ oraz
$\fint_E h\,\mathrm d\lambda_d=|E|^{-1}
\int_E h\,\mathrm d\lambda_d$ dla zbiorów mierzalnych o dodatniej skończonej mierze;
$|E|=\lambda_d(E)$, a $h_E$ oznacza średnią funkcji
całkowalnej na $E$. Porównania wartości funkcji Sobolewa rozumiemy prawie
wszędzie, natomiast słabe nierówności różniczkowe
definiujemy za pomocą funkcji testowych.

\section{Założenia, słabe nierówności i skalowanie}
\label{sec:pde17-assumptions}
\begin{definicja}[Jednostajna eliptyczność i rozwiązania słabe]
Niech $d\ge2$, $\Omega\subset\R^d$ będzie obszarem, a
$A\colon\Omega\to\R^{d\times d}$ mierzalną macierzą, taką że dla
prawie każdego $x\in\Omega$, $\xi,\zeta\in\R^d$ zachodzą
\begin{equation}\label{eq:pde17-ellipticity}
 \xi^{\mathsf T}A(x)\xi\ge\lambda|\xi|^2,
 \qquad |A(x)\xi\cdot\zeta|\le\Lambda|\xi||\zeta|,
 \qquad 0<\lambda\le\Lambda<\infty.
\end{equation}
Nie wymagamy symetrii $A$. Rzeczywistą funkcję
$u\in H^1_{\rm loc}(\Omega)$
nazywamy słabym rozwiązaniem równania
$-\operatorname{div}(A\nabla u)=0$, jeżeli
\begin{equation}\label{eq:pde17-weak}
 \int_\Omega A\nabla u\cdot\nabla\varphi\,\mathrm d\lambda_d=0
 \quad\text{dla każdej }\varphi\in C_c^\infty(\Omega).
\end{equation}
Słabym podrozwiązaniem (odpowiednio nadrozwiązaniem)
nazywamy rzeczywistą funkcję $u\in H^1_{\rm loc}(\Omega)$,
dla której lewa strona~\eqref{eq:pde17-weak} jest $\le0$
(odpowiednio $\ge0$) dla każdej nieujemnej funkcji testowej.
\end{definicja}

\begin{uwaga}[Dopuszczalne testy]
Niech $\varphi\in H^1(\Omega)$ ma zwarty nośnik $K\Subset\Omega$.
Wybieramy zbiór otwarty $V$ taki, że $K\subset V\Subset\Omega$.
Przedłużenie $\varphi$ przez zero poza $\Omega$ należy do
$H^1(\R^d)$: można najpierw pomnożyć $\varphi$ przez gładką
funkcję odcięcia równą $1$ w otoczeniu $K$ i o nośniku w $V$,
a następnie przedłużyć iloczyn przez zero.
Sploty tego przedłużenia z nieujemnym gładkim jądrem
aproksymującym jedynkę, o dostatecznie małym nośniku, należą
do $C_c^\infty(V)$ i zbiegają do $\varphi$ w $H^1(V)$.
Jeśli $\varphi\ge0$ prawie wszędzie, także sploty są nieujemne.
Ponieważ
\[
 \left|\int_V A\nabla u\cdot\nabla(\varphi_j-\varphi)\,
 \mathrm d\lambda_d\right|
 \le\Lambda\|\nabla u\|_{L^2(V)}
       \|\nabla(\varphi_j-\varphi)\|_{L^2(V)}\longrightarrow0,
\]
równość oraz odpowiednie słabe nierówności przechodzą na takie
testy, z zachowaniem wymaganego znaku.
Dopuszczalność obcięć i potęg sprawdzamy za każdym razem za
pomocą reguły łańcuchowej dla funkcji Sobolewa i zwartego
nośnika funkcji odcięcia. Nie zakładamy gładkości $u$
ani nie różniczkujemy współczynników $A$.
\end{uwaga}

Jeżeli $u$ rozwiązuje równanie w $B_{2R}(x_0)$, funkcja
$v(y)=u(x_0+Ry)$ rozwiązuje je w $B_2(0)$ z macierzą
$A_R(y)=A(x_0+Ry)$ i \emph{tymi samymi} stałymi $\lambda,\Lambda$.
Istotnie, dla $\psi\in C_c^\infty(B_2)$ testujemy równanie
funkcją $\varphi(x)=\psi((x-x_0)/R)$. Z reguły łańcuchowej
dla słabych pochodnych i podstawienia $x=x_0+Ry$ otrzymujemy
\[
 \int_{B_2} A_R\nabla v\cdot\nabla\psi\,\mathrm d\lambda_d
 =R^{2-d}\int_{B_{2R}(x_0)}
 A\nabla u\cdot\nabla\varphi\,\mathrm d\lambda_d=0.
\]
Dodatni czynnik $R^{2-d}$ zachowuje także znak słabej
nierówności. Dzieląc $A_R$ przez $\lambda$, normalizujemy
dolną stałą eliptyczności do $1$, a górną do $\Lambda/\lambda$.
Stałe w znormalizowanych lokalnych oszacowaniach zależą
więc od $d$ i $\Lambda/\lambda$, a odpowiednie potęgi
promienia wynikają z powrotu do pierwotnej skali.

\section{Nierówność Caccioppolego dla obcięć De Giorgiego}
\label{sec:pde17-caccioppoli}
\begin{lemat}[Energia części dodatniej]
\label{lem:pde17-caccioppoli}
Niech $u\in H^1(B_R)$ będzie podrozwiązaniem w $B_R\subset\Omega$.
Dla $k\in\R$ i $\eta\in C_c^1(B_R)$, $0\le\eta\le1$, zachodzi
\begin{equation}\label{eq:pde17-cacc}
 \int_{B_R}\eta^2|\nabla(u-k)_+|^2\,\mathrm d\lambda_d
 \le \frac{4\Lambda^2}{\lambda^2}
 \int_{B_R}(u-k)_+^2|\nabla\eta|^2\,\mathrm d\lambda_d.
\end{equation}
Analogicznie, dla nadrozwiązania oszacowanie dotyczy $(k-u)_+$.
\end{lemat}
\begin{proof}
Oznaczmy $w=(u-k)_+$. Funkcja $w$ należy do $H^1$;
$\nabla w=\chi_{\{u>k\}}\nabla u$ prawie wszędzie,
co wynika z wniosku~\ref{cor:pde9-truncation}.
Testujemy nierówność podrozwiązania funkcją $\varphi=\eta^2w\ge0$.
Ponieważ $\nabla\varphi=\eta^2\nabla w+2\eta w\nabla\eta$
prawie wszędzie, a oba składniki znikają na $\{u\le k\}$,
dostajemy
\[
 \lambda\int_{B_R}\eta^2|\nabla w|^2\,\mathrm d\lambda_d
 \le 2\Lambda\int_{B_R}\eta w|\nabla w||\nabla\eta|\,\mathrm d\lambda_d.
\]
Nierówność Younga
$2\Lambda ab\le(\lambda/2)a^2+(2\Lambda^2/\lambda)b^2$
daje tezę. Jeżeli $u$ jest nadrozwiązaniem, to $-u$
jest podrozwiązaniem, gdyż forma jest liniowa względem
pierwszego argumentu. Stosujemy już udowodnione oszacowanie
do $-u$ i poziomu $-k$, zauważając, że
$(-u-(-k))_+=(k-u)_+$.
\end{proof}

\begin{uwaga}[Kontrast z metodą ilorazów różnicowych]
W poprzednich rozdziałach do regularności $H^2$ wykorzystywaliśmy
pochodne $A$. W~\eqref{eq:pde17-cacc} wystarcza informacja o
$A\in L^\infty$ i dodatniości części symetrycznej macierzy.
To właśnie pozwala prowadzić argument dla współczynników mierzalnych.
\end{uwaga}

\section{Lemat iteracyjny De Giorgiego i lokalne ograniczenie}
\label{sec:pde17-degiorgi}
\begin{lemat}[Zbieżność iteracji superliniowej]
\label{lem:pde17-iteration}
Niech $Y_n\ge0$, $a>1$, $b>1$, $\delta>0$ oraz
$Y_{n+1}\le a b^n Y_n^{1+\delta}$ dla $n\ge0$.
Jeżeli $Y_0\le a^{-1/\delta}b^{-1/\delta^2}$,
to $Y_n\to0$.
\end{lemat}
\begin{proof}
Niech $Z_n=a^{1/\delta}b^{n/\delta+1/\delta^2}Y_n$.
Podstawienie do rekurencji daje
\[
 Z_{n+1}\le
 b^{(n+1)/\delta+1/\delta^2+n
 -(1+\delta)(n/\delta+1/\delta^2)}
 Z_n^{1+\delta}=Z_n^{1+\delta},
\]
gdyż potęgi $a$ skracają się, a wykładnik $b$ jest zerowy.
Z założenia $0\le Z_0\le1$, więc przez indukcję
$0\le Z_n\le1$ dla wszystkich $n$, a zatem
$Y_n\le a^{-1/\delta}b^{-n/\delta-1/\delta^2}\to0$.
\end{proof}

\begin{twierdzenie}[Oszacowanie $L^2\to L^\infty$ podrozwiązania]
\label{thm:pde17-l2linfty}
Niech $B_{2R}(x_0)\Subset\Omega$ i $u\in H^1(B_{2R}(x_0))$
będzie podrozwiązaniem równania z~\eqref{eq:pde17-ellipticity}.
Wtedy
\begin{equation}\label{eq:pde17-l2linfty}
 \operatorname*{ess\,sup}_{B_R(x_0)}u_+
 \le C\left(\fint_{B_{2R}(x_0)}u_+^2\,\mathrm d\lambda_d\right)^{1/2},
 \qquad C=C(d,\Lambda/\lambda).
\end{equation}
\end{twierdzenie}
\begin{proof}
Skalowanie pozwala przyjąć $R=1$, $x_0=0$. Ustalmy
$K>0$ i ciągi
$r_n=1+2^{-n}$, $k_n=K(1-2^{-n-1})$, $n\ge0$.
Niech $w_n=(u-k_n)_+$ oraz
$Y_n=\int_{B_{r_n}}w_n^2\,\mathrm d\lambda_d$.
Wybieramy $\eta_n=1$ na $B_{r_{n+1}}$, o nośniku w $B_{r_n}$,
z $|\nabla\eta_n|\le C2^n$.
Przedłużamy $\eta_n w_n$ przez zero do $B_2$.
Z nierówności Sobolewa (twierdzenie~\ref{thm:pde10-Sobp})
dla $d>2$, a dla $d=2$ z tej samej nierówności dla $p=4/3$
i nierówności Höldera na $B_2$, otrzymujemy dla odpowiedniego
$\chi>1$
\[
 \|\eta_n w_n\|_{L^{2\chi}(B_2)}^2
 \le C_d\|\nabla(\eta_n w_n)\|_{L^2(B_2)}^2.
\]
Stosujemy te oszacowania do przedłużeń przez zero, korzystając
z gęstości $C_c^\infty(B_2)$ w $H_0^1(B_2)$.
Z reguły Leibniza i lematu~\ref{lem:pde17-caccioppoli}
prawa strona jest nie większa od
$C\int_{B_{r_n}}w_n^2|\nabla\eta_n|^2\,\mathrm d\lambda_d$,
a więc
\begin{equation}\label{eq:pde17-sobstep}
 \left(\int_{B_{r_n}}|\eta_nw_n|^{2\chi}\,\mathrm d\lambda_d
 \right)^{1/\chi}\le C4^nY_n.
\end{equation}
Dla $d>2$ bierzemy $\chi=d/(d-2)$, a dla $d=2$
możemy wybrać $\chi=2$ (zanurzenie $H_0^1(B_2)\hookrightarrow L^4$).
Na zbiorze $E_{n+1}=B_{r_{n+1}}\cap\{u>k_{n+1}\}$ mamy
$w_n\ge k_{n+1}-k_n=K2^{-n-2}$, zatem
$|E_{n+1}|\le C4^nK^{-2}Y_n$. Nierówność Höldera i
\eqref{eq:pde17-sobstep} dają
\[
 Y_{n+1}\le
 \left(\int_{B_{r_n}}|\eta_nw_n|^{2\chi}\,\mathrm d\lambda_d\right)^{1/\chi}
 |E_{n+1}|^{1-1/\chi}
 \le Cb^nK^{-2\delta}Y_n^{1+\delta},
 \quad\delta=1-1/\chi>0,
\]
przy pewnym $b=b(d)>1$. Po podzieleniu przez $K^2$
otrzymujemy rekurencję dla $Z_n=Y_n/K^2$ postaci
z lematu~\ref{lem:pde17-iteration}. Jeżeli $\|u_+\|_{L^2(B_2)}=0$, teza jest natychmiastowa.
W przeciwnym razie bierzemy
$K=C_0\|u_+\|_{L^2(B_2)}>0$.
Ponieważ $Y_0\le\|u_+\|_{L^2(B_2)}^2$, wybór dostatecznie
dużego $C_0$ zapewnia próg małości dla $Z_0$ z lematu
~\ref{lem:pde17-iteration}. Otrzymujemy $Y_n\to0$.
Na $B_1\subset B_{r_n}$ mamy $(u-K)_+\le w_n$, więc
$\int_{B_1}(u-K)_+^2\,\mathrm d\lambda_d\le Y_n\to0$,
czyli $u\le K$ prawie wszędzie na $B_1$.
Dla $v(y)=u(x_0+Ry)$ zachodzi
\[
 \|v_+\|_{L^2(B_2)}
 =|B_2|^{1/2}
 \left(\fint_{B_{2R}(x_0)}u_+^2\,\mathrm d\lambda_d\right)^{1/2}.
\]
Wchłaniając $|B_2|^{1/2}$ do stałej zależnej od $d$,
otrzymujemy~\eqref{eq:pde17-l2linfty}.
\end{proof}

\begin{wniosek}[Lokalna ograniczoność rozwiązania]
Jeżeli $u$ jest rozwiązaniem słabym w $B_{2R}$, to
$u$ i $-u$ są podrozwiązaniami. Stąd
$\|u\|_{L^\infty(B_R)}\le C(\fint_{B_{2R}}|u|^2\,\mathrm d\lambda_d)^{1/2}$.
W szczególności każde rozwiązanie słabe jest lokalnie
istotnie ograniczone.
\end{wniosek}

\section{Energia potęg i mechanizm iteracji Mosera}
\label{sec:pde17-moser}
W dalszej części zakładamy $v\ge\eps>0$. Dla nieujemnego
rozwiązania $u$ otrzymujemy to, biorąc $v=u+\eps$,
a po udowodnieniu oszacowań przechodzimy do granicy
$\eps\downarrow0$. Takie przesunięcie jest rozwiązaniem,
ponieważ równanie nie zawiera wyrazu zerowego rzędu.
Ponadto z lokalnej ograniczoności rozwiązania otrzymujemy
$v\in L^\infty_{\rm loc}$. Na zwartym nośniku każdego testu
wartości $v$ należą prawie wszędzie do pewnego przedziału
$[\eps,M]$. Funkcję potęgową można przedłużyć z tego przedziału
do funkcji klasy $C^1$ o ograniczonej pochodnej na $\R$
(np.\ przez liniową kontynuację z obu końców przedziału).
Stosujemy wówczas twierdzenie~\ref{thm:pde9-chain}
na otoczeniu nośnika testu, na którym $v$ jest ograniczona. Stałe w uzyskanych oszacowaniach
nie będą zależeć od $\eps$ ani od $M$.

\begin{lemat}[Oszacowanie dla potęg dodatnich i ujemnych]
\label{lem:pde17-powers}
Niech $v\in H^1(B_R)\cap L^\infty(B_R)$, $v\ge\eps>0$,
będzie słabym rozwiązaniem, $\beta\in(-\infty,0)\cup(1,\infty)$,
a $\eta\in C_c^1(B_R)$, $0\le\eta\le1$.
Wówczas
\begin{equation}\label{eq:pde17-powerenergy}
 \int_{B_R}|\nabla(\eta v^{\beta/2})|^2\,\mathrm d\lambda_d
 \le C\left(1+\frac{\beta^2}{(\beta-1)^2}\right)
 \int_{B_R}v^\beta|\nabla\eta|^2\,\mathrm d\lambda_d,
\end{equation}
gdzie $C$ zależy tylko od $\Lambda/\lambda$.
W szczególności stała jest ograniczona jednostajnie dla
$\beta\ge2$ lub $\beta<0$.
\end{lemat}
\begin{proof}
Testujemy równanie funkcją $\eta^2v^{\beta-1}$.
Jest to dopuszczalne, ponieważ $\eps\le v\le M$ na nośniku
$\eta$; mapy potęgowe mają tam ograniczone pochodne.
Reguła łańcuchowa i reguła Leibniza dają
\[
 0=(\beta-1)\int_{B_R}\eta^2v^{\beta-2}
 A\nabla v\cdot\nabla v\,\mathrm d\lambda_d
 +2\int_{B_R}\eta v^{\beta-1}
 A\nabla v\cdot\nabla\eta\,\mathrm d\lambda_d.
\]
Pierwsza całka bez czynnika $\beta-1$ jest nieujemna.
Biorąc wartości bezwzględne po przeniesieniu drugiego
składnika na prawą stronę, otrzymujemy niezależnie od znaku
$\beta-1$
\[
 |\beta-1|\lambda\int_{B_R}\eta^2v^{\beta-2}|\nabla v|^2
 \,\mathrm d\lambda_d
 \le2\Lambda\int_{B_R}\eta v^{\beta-1}|\nabla v||\nabla\eta|
 \,\mathrm d\lambda_d.
\]
Stosujemy nierówność Younga do iloczynu
$a=\eta v^{(\beta-2)/2}|\nabla v|$ oraz
$b=v^{\beta/2}|\nabla\eta|$:
\[
 2\Lambda ab\le\frac{\lambda|\beta-1|}{2}a^2+
 \frac{2\Lambda^2}{\lambda|\beta-1|}b^2.
\]
Po absorpcji pierwszego składnika otrzymujemy
\[
 \int_{B_R}a^2\,\mathrm d\lambda_d
 \le\frac{4(\Lambda/\lambda)^2}{(\beta-1)^2}
 \int_{B_R}b^2\,\mathrm d\lambda_d.
\]
Ponieważ $|\nabla(v^{\beta/2})|^2
=(\beta^2/4)v^{\beta-2}|\nabla v|^2$,
zaś $|\nabla(\eta v^{\beta/2})|^2$ jest co najwyżej
$2\eta^2|\nabla(v^{\beta/2})|^2+2v^\beta|\nabla\eta|^2$,
otrzymujemy~\eqref{eq:pde17-powerenergy}.
\end{proof}

\begin{lemat}[Iteracja $L^q\to L^\infty$ i potęg odwrotnych]
\label{lem:pde17-moser-iteration}
Niech $B_{2R}\Subset\Omega$, $v$ będzie rozwiązaniem słabym
w $\Omega$ oraz $v\ge\eps>0$ na $B_{2R}$, a $0<r<s\le2R$.
Dla każdego $q>0$ istnieje $C=C(d,\Lambda/\lambda,q)$ takie, że
\begin{align}
 \operatorname*{ess\,sup}_{B_r}v
 &\le C(s-r)^{-d/q}\|v\|_{L^q(B_s)},\label{eq:pde17-moser-up}\\
 \operatorname*{ess\,sup}_{B_r}v^{-1}
 &\le C(s-r)^{-d/q}\|v^{-1}\|_{L^q(B_s)}.
 \label{eq:pde17-moser-inv}
\end{align}
Dla $0<q<1$ zapis $\|w\|_{L^q(E)}$ oznacza wielkość
$(\int_E|w|^q\,\mathrm d\lambda_d)^{1/q}$, nie normę.
Lokalna ograniczoność $v$ wynika z twierdzenia
~\ref{thm:pde17-l2linfty}; dla $u\ge0$ stosujemy
nierówności do $v=u+\eps$.
\end{lemat}
\begin{proof}
Najpierw dowodzimy oszacowania na parach kul $B_1\subset B_2$
dla $q\ge2$. Przy $d>2$ bierzemy $\chi=d/(d-2)$, a przy
$d=2$ wybieramy $\chi=2$. Na ustalonej kuli $B_2$ stosujemy to samo oszacowanie
Sobolewa dla $H_0^1(B_2)$ co w dowodzie twierdzenia
~\ref{thm:pde17-l2linfty}:
$\|h\|_{L^{2\chi}(B_2)}^2\le C_d\|\nabla h\|_{L^2(B_2)}^2$.
Dla $1\le\rho<\sigma\le2$ wybieramy odcięcie równe $1$
na $B_\rho$, z nośnikiem w $B_\sigma$ i pochodną ograniczoną
przez $C/(\sigma-\rho)$. Z~\eqref{eq:pde17-powerenergy},
dla $|\beta|\ge2$, wynika
\[
 \left(\int_{B_\rho}v^{\beta\chi}\,\mathrm d\lambda_d\right)^{1/\chi}
 \le\frac{C}{(\sigma-\rho)^2}
 \int_{B_\sigma}v^\beta\,\mathrm d\lambda_d.
\]
Niech $w=v$ dla $\beta>0$ i $w=v^{-1}$ dla $\beta<0$.
Po podniesieniu do dodatniej potęgi $1/|\beta|$ otrzymujemy
\[
 \|w\|_{L^{|\beta|\chi}(B_\rho)}
 \le\left(\frac{C}{(\sigma-\rho)^2}\right)^{1/|\beta|}
       \|w\|_{L^{|\beta|}(B_\sigma)}.
\]
Bierzemy $p_k=q\chi^k$, $\rho_k=1+2^{-k}$.
Kolejne zastosowania z $|\beta|=p_k$ dają
\[
 \|w\|_{L^{p_n}(B_{\rho_n})}
 \le \prod_{k=0}^{n-1}(C4^{k+1})^{1/p_k}
        \|w\|_{L^q(B_2)}.
\]
Iloczyn jest ograniczony niezależnie od $n$, ponieważ
$\sum_{k\ge0}(1+k)/\chi^k<\infty$.
Na $B_1$ norma $L^{p_n}$ jest nie większa od normy na
$B_{\rho_n}$. Dla każdego $a>0$, dla którego zbiór
$E_a=B_1\cap\{w>a\}$ ma dodatnią miarę, mamy
$a|E_a|^{1/p_n}\le\|w\|_{L^{p_n}(B_1)}
\le C\|w\|_{L^q(B_2)}$.
Ponieważ $|E_a|^{1/p_n}\to1$, dostajemy
$a\le C\|w\|_{L^q(B_2)}$. Biorąc supremum po takich
$a$ (przypadek $w=0$ prawie wszędzie jest natychmiastowy),
otrzymujemy $\|w\|_{L^\infty(B_1)}\le C\|w\|_{L^q(B_2)}$.

Przeskalowanie pary $B_t(y)\subset B_{2t}(y)$ do powyższych
kul daje
\[
 \|w\|_{L^\infty(B_t(y))}
 \le C t^{-d/q}\|w\|_{L^q(B_{2t}(y))}.
\]
W szczególności wykładnik $d/q$ jest poprawny także dla $d=2$:
skalujemy ustaloną parę kul, a nie stosujemy jednorodnej
nierówności Sobolewa z wykładnikiem $4$ na dowolnej skali.
Dla $0<r<s$ kładziemy $t=(s-r)/2$. Dla każdego $y\in B_r$
mamy $B_{2t}(y)\subset B_s$. Przeliczalne pokrycie $B_r$
kulami $B_t(y)$ (np.\ o środkach z przeliczalnego zbioru
gęstego w $B_r$) pozwala połączyć oszacowania prawie wszędzie:
\[
 \|w\|_{L^\infty(B_r)}
 \le C(s-r)^{-d/q}\|w\|_{L^q(B_s)}.
\]

Dla $0<q<2$ z punktowej nierówności
$|w|^2\le\|w\|_{L^\infty(B_\sigma)}^{2-q}|w|^q$
po scałkowaniu wynika
$\|w\|_{L^2(B_\sigma)}\le
\|w\|_{L^\infty(B_\sigma)}^{1-q/2}
\|w\|_{L^q(B_\sigma)}^{q/2}$,
gdzie $w=v$ lub $w=v^{-1}$. Aby uzasadnić absorpcję,
ustalmy $r_j=r+(s-r)(1-2^{-j})$, $j\ge0$, i
$M_j=\|w\|_{L^\infty(B_{r_j})}$ oraz
$Q=\|w\|_{L^q(B_s)}$. Stosując oszacowanie $L^2\to
L^\infty$ pomiędzy $B_{r_j}$ a $B_{r_{j+1}}$
otrzymujemy
\[
 M_j\le C2^{jd/2}(s-r)^{-d/2}M_{j+1}^{1-q/2}Q^{q/2}.
\]
Dla dowolnego $0<\delta<1$ nierówność Younga z wykładnikami
$2/(2-q)$ i $2/q$ daje
\[
 M_j\le\delta M_{j+1}
 +C_{d,\Lambda/\lambda,q,\delta}2^{jd/q}(s-r)^{-d/q}Q.
\]
Wybieramy $\delta$ tak małe, aby $\delta2^{d/q}<1$,
i iterujemy nierówność $n$ razy:
\[
 M_0\le\delta^n M_n+
 C(s-r)^{-d/q}Q\sum_{j=0}^{n-1}(\delta2^{d/q})^j.
\]
Szereg geometryczny jest zbieżny. Pozostałość
$\delta^nM_n$ dąży do zera: dzięki lokalnej ograniczoności
$w$ jest ograniczone na $B_s\Subset\Omega$. Uzyskujemy
$M_0\le C(s-r)^{-d/q}Q$, czyli
\eqref{eq:pde17-moser-up}--\eqref{eq:pde17-moser-inv}.
Dowodzi to oszacowań także dla $0<q<2$.
\end{proof}


\begin{uwaga}[Jawny zakres stosowania potęg ujemnych]
\label{rem:pde17-negpower}
W oszacowaniu~\eqref{eq:pde17-moser-inv} nie wolno od razu
podstawić $u^{-1}$, jeśli nie wykazano dodatniej dolnej granicy
rozwiązania. Dla $u\ge0$ pracujemy z $v=u+\eps$;
wówczas $v^{-1}\in L^\infty$ na każdej kuli, a reguła
łańcuchowa jest dopuszczalna. Najpierw dowodzimy obu
oszacowań dla każdego $\eps>0$, przy tych samych stałych,
a dopiero na końcu przechodzimy do granicy monotonnej.
Nie otrzymujemy w ten sposób osobnego oszacowania
$\|u^{-1}\|_{L^q}$, które mogłoby być nieskończone.
To właśnie połączenie oszacowania dla $v^{-1}$
z mostem~\eqref{eq:pde17-bridge} usuwa tę trudność.
\end{uwaga}

\section{Logarytm rozwiązania i wykładnicza całkowalność}
\label{sec:pde17-log}
\begin{lemat}[Energia logarytmu]
\label{lem:pde17-log}
Niech $v\ge\eps>0$ będzie słabym dodatnim nadrozwiązaniem
w $B_{2r}(z)$. Wówczas
\begin{equation}\label{eq:pde17-log}
 \int_{B_r(z)}|\nabla\log v|^2\,\mathrm d\lambda_d
 \le C r^{d-2},\qquad C=C(d,\Lambda/\lambda).
\end{equation}
Dla każdej kuli $B_s(y)\subset B_r(z)$ otrzymujemy w konsekwencji
$\fint_{B_s(y)}|\log v-(\log v)_{B_s(y)}|\,\mathrm d\lambda_d
\le C(d,\Lambda/\lambda)$, o ile równanie obowiązuje
na kuli o podwójnym promieniu.
\end{lemat}
\begin{proof}
Niech $\eta\in C_c^1(B_{2r}(z))$, $0\le\eta\le1$.
Funkcje $t\mapsto t^{-1}$ i $t\mapsto\log t$ na
$[\eps,\infty)$ przedłużamy liniowo poniżej $\eps$,
otrzymując funkcje klasy $C^1$ o ograniczonych pochodnych.
Z twierdzenia~\ref{thm:pde9-chain}, stosowanego lokalnie,
wynika $\log v\in H^1_{\rm loc}$, z
$\nabla\log v=v^{-1}\nabla v$ prawie wszędzie, oraz
$\varphi=\eta^2v^{-1}$ jest dopuszczalnym nieujemnym
testem o zwartym nośniku. Z nierówności nadrozwiązania
uzyskujemy
\[
 \int_{B_{2r}(z)}\eta^2v^{-2}A\nabla v\cdot\nabla v\,\mathrm d\lambda_d
 \le 2\int_{B_{2r}(z)}\eta v^{-1}A\nabla v\cdot\nabla\eta\,\mathrm d\lambda_d.
\]
Eliptyczność, ograniczenie $A$ i nierówność Younga
z $a=\eta|\nabla\log v|$, $b=|\nabla\eta|$ dają
\[
 \int_{B_{2r}(z)}\eta^2|\nabla\log v|^2\,\mathrm d\lambda_d
 \le 4(\Lambda/\lambda)^2
 \int_{B_{2r}(z)}|\nabla\eta|^2\,\mathrm d\lambda_d.
\]
Wybieramy $\eta=1$ na $B_r(z)$ z nośnikiem w $B_{2r}(z)$,
$|\nabla\eta|\le C_d/r$. Całka po prawej stronie jest
nie większa niż $C_d r^{d-2}$.
Dla $B=B_s(y)$, $g=\log v$ i
$g_B=\fint_B g\,\mathrm d\lambda_d$ korzystamy z nierówności
Poincarégo (wniosek~\ref{cor:pde10-Poincare}, zastosowany
na kuli jednostkowej i przeskalowany) i nierówności Höldera:
\[
 \fint_B |g-g_B|\,\mathrm d\lambda_d
 \le C_d s\left(\fint_B|\nabla g|^2\,\mathrm d\lambda_d\right)^{1/2}
 \le C(d,\Lambda/\lambda).
\]
Ostatnia nierówność wynika z oszacowania energii na $B_s(y)$,
otrzymanego przy użyciu odcięcia na $B_{2s}(y)$.
\end{proof}

Potrzebujemy teraz jednego lematu analizy harmonicznej.
Nie traktujemy go jako automatycznej własności funkcji $H^1$:
przydatna jest \emph{jednostajna} kontrola średnich oscylacji
na wszystkich podkulach.

\begin{lemat}[Lemat Johna--Nirenberga; lokalna postać]
\label{lem:pde17-jn}
Niech $B=B_R(z)\subset\R^d$, a
$g\in L^1(B_{3R/2}(z))$. Zakładamy, że
$\sup_{B'\subset B_{3R/2}(z)}\fint_{B'}
|g-g_{B'}|\,\mathrm d\lambda_d\le M$,
gdzie supremum biegnie po kulach zawartych w większej kuli.
Jeśli $M>0$, istnieją $c_d,C_d>0$, dla których
\begin{equation}\label{eq:pde17-jn}
 \fint_B\exp\left(\frac{c_d}{M}|g-g_B|\right)
 \,\mathrm d\lambda_d\le C_d
\end{equation}
(gdy $M=0$, funkcja jest stała prawie wszędzie).
\end{lemat}
\begin{proof}[Dowód metodą kostek zatrzymania]
Gdy $M=0$, możemy zastosować założenie do samej kuli $B$;
wtedy $\int_B|g-g_B|\,\mathrm d\lambda_d=0$, więc $g=g_B$
prawie wszędzie na $B$. Dalej zakładamy $M>0$.
Najpierw ustalmy kostkę $Q$ i załóżmy, że średnia oscylacja
jest ograniczona przez $M$ na wszystkich jej podkostkach.
Rozkładamy $Q$ na kostki dyadyczne, pomijając ich brzegi,
których przeliczalna suma ma miarę zero.
Wybieramy $L=2^{d+2}M$. Ponieważ średnia oscylacja
na całej $Q$ jest nie większa niż $M<L$, każda kostka
przekraczająca próg zawiera się w maksymalnej takiej
kostce (wystarczy przejść po skończonym ciągu jej przodków).
Maksymalne kostki dyadyczne
$Q_j\subset Q$, dla których
$\fint_{Q_j}|g-g_Q|>L$, są rozłączne. Ponieważ
$\int_Q|g-g_Q|\le M|Q|$, łączna miara tych kostek
jest nie większa od $(M/L)|Q|\le 2^{-d-2}|Q|$.
Maksymalność i porównanie z rodzicem dają
$|g_{Q_j}-g_Q|\le2^dL$. Poza ich sumą średnia $|g-g_Q|$ na każdej kostce
dyadycznej zawierającej dany punkt nie przekracza $L$.
W punkcie Lebesgue'a funkcji $|g-g_Q|$ średnie te dążą
do jej wartości: kostka boku $b$ zawierająca ten punkt
mieści się w kuli o środku w tym punkcie i promieniu
$\sqrt d\,b$, a iloraz objętości kuli i kostki zależy
tylko od $d$. Z twierdzenia o różniczkowaniu całek
Lebesgue'a wynika więc $|g-g_Q|\le L$ prawie wszędzie
poza sumą wybranych kostek. W każdej złej kostce powtarzamy procedurę
z $g_{Q_j}$ zamiast $g_Q$. Oznaczmy przez $G_k$ sumę
złych kostek $k$-tego pokolenia, przyjmując $G_0=Q$.
W każdej kostce poprzedniego pokolenia suma miar
wybranych potomków stanowi co najwyżej część $M/L$
jej miary. Rozłączność kostek w danym pokoleniu i indukcja
dają $|G_k|\le (M/L)^k|Q|=2^{-k(d+2)}|Q|$.
Dla prawie każdego $x\in Q\setminus G_k$ łańcuch
wybranych kostek zawierających $x$ kończy się
po $j<k$ krokach. W ostatniej kostce $P$ zachodzi
$|g(x)-g_P|\le L$, a różnica średnich między każdą
parą kolejnych kostek w tym łańcuchu jest ograniczona
przez $2^dL$. Nierówność trójkąta daje więc
$|g(x)-g_Q|\le(1+k2^d)L$. Dla $k=0$ oszacowanie
na dopełnieniu jest puste.
Dla $t\ge L$ wybieramy największe $k\ge0$ takie, że
$(1+k2^d)L\le t$. Z oszacowania miary $k$-tego pokolenia,
uwzględniając $L=2^{d+2}M$, wynika
$|\{x\in Q:|g(x)-g_Q|>t\}|\le C|Q|e^{-ct/M}$.
Po zwiększeniu $C$ oszacowanie obowiązuje także dla $t<L$.
Dla $X=|g-g_Q|/M$ i $0<c'<c$ z tożsamości
$e^{c'X}=1+\int_0^X c'e^{c't}\,\mathrm dt$ oraz
twierdzenia Tonellego otrzymujemy
\[
 \fint_Q e^{c'X}\,\mathrm d\lambda_d
 =1+\frac{1}{|Q|}\int_0^\infty c'e^{c't}
      |\{x\in Q:X(x)>t\}|\,\mathrm dt
 \le 1+\frac{Cc'}{c-c'}.
\]
Przejście do kuli wymaga jeszcze uzasadnienia. Ustalmy
$B=B_R(z)$ i wybierzmy siatkę kostek o boku
$\ell=R/(100\sqrt d)$. Niech $h=\ell/2$. Wybieramy wszystkie otwarte kostki
o boku $\ell$, o środkach $y\in z+h\mathbb Z^d$ spełniających
$|y-z|\le R+\sqrt d\,h/2$. Dla każdego punktu
$\overline B$ zaokrąglenie współrzędnych do najbliższego
punktu siatki daje taki środek, z odległością w każdej
współrzędnej co najwyżej $h/2<\ell/2$. Kostki pokrywają
więc $\overline B$. Z każdego środka można dojść do $z$
krokami długości $h$ wzdłuż osi, zmniejszając kolejno
wartości bezwzględne współrzędnych względem $z$.
Odległość od $z$ nie rośnie, zatem wszystkie środki
pośrednie należą do pokrycia. Kostki odpowiadające
kolejnym krokom mają przecięcie objętości $\ell^d/2$.
Każde trzykrotne powiększenie kostki mieści się
w $B_{3R/2}(z)$, gdyż
$R+\sqrt d\,h/2+3\sqrt d\,\ell/2<3R/2$.
Liczba kostek $N$ jest ograniczona przez stałą zależną
tylko od $d$, ponieważ $R/h=200\sqrt d$.

Na każdej podkostce $Q'\subset Q_j$ średnia oscylacja
$g$ nie przekracza $C_dM$. Rzeczywiście, otaczamy $Q'$
kulą $B'$ o środku w środku $Q'$ i promieniu
$\sqrt d$ razy bok $Q'$ (można też wybrać promień
$\sqrt d/2$ razy bok $Q'$ dla kostki otwartej).
Jeśli $b$ jest bokiem $Q'$, to odległość środka
$Q'$ od środka $Q_j$ jest nie większa niż
$\sqrt d(\ell-b)/2$. Zatem $B'$ mieści się w kuli
o środku w środku $Q_j$ i promieniu
$\sqrt d(\ell+b)/2\le\sqrt d\,\ell$.
Ta kula wraz z $Q_j$ znajduje się w $B_{3R/2}(z)$,
na mocy ograniczenia odległości środków od $z$. Z nierówności
$\fint_{Q'}|g-g_{Q'}|\le
2\fint_{Q'}|g-g_{B'}|\le
2(|B'|/|Q'|)\fint_{B'}|g-g_{B'}|$
dostajemy żądane ograniczenie, ponieważ iloraz objętości
zależy tylko od $d$. Dyadyczny dowód dla kostek
zastosowany do każdego $Q_j$ zapewnia zatem
\[
 \fint_{Q_j}\exp\bigl(c_d|g-g_{Q_j}|/M\bigr)
 \,\mathrm d\lambda_d\le C_d.
\]
Jeżeli sąsiednie kostki $Q_i,Q_j$ mają część wspólną
o objętości co najmniej $c_d|Q_i|$, to
$|g_{Q_i}-g_{Q_j}|\le C_dM$: wystarczy całkować
$|g_{Q_i}-g_{Q_j}|\le
|g-g_{Q_i}|+|g-g_{Q_j}|$ po ich przecięciu,
a oba składniki oszacować przez oscylacje na kostkach.
Każdą parę kostek pokrycia łączy łańcuch co najwyżej
$N-1$ takich sąsiadów, stąd
$|g_{Q_j}-g_{Q_1}|\le C_dM$ dla wszystkich $j$.
Z nierówności trójkąta mamy dla $x\in Q_j$
$e^{c_d|g(x)-g_{Q_1}|/M}
\le e^{c_d C_d}e^{c_d|g(x)-g_{Q_j}|/M}$.
Ponieważ $B\subset\bigcup_{j=1}^N Q_j$, całkę po $B$
ograniczamy sumą całek po kostkach. Po podzieleniu
przez $|B|$ otrzymujemy
\[
 \fint_B e^{c_d|g-g_{Q_1}|/M}\,\mathrm d\lambda_d
 \le C_d\frac{N\ell^d}{|B|}\le C'_d,
\]
gdzie ostatnia stała zależy wyłącznie od $d$,
na mocy konstrukcji pokrycia.
Z nierówności Jensena lub bezpośrednio z oszacowania
pierwszego momentu wynika
$|g_B-g_{Q_1}|\le C_dM$. Zastępując $g_{Q_1}$
przez $g_B$, kończymy dowód~\eqref{eq:pde17-jn}.
\end{proof}

\begin{wniosek}[Most między potęgami dodatnimi i ujemnymi]
\label{cor:pde17-bridge}
Jeżeli $v\ge\eps>0$ jest nadrozwiązaniem w $B_{4R}(x_0)$,
to istnieje $q_0=q_0(d,\Lambda/\lambda)>0$, dla którego
\begin{equation}\label{eq:pde17-bridge}
 \left(\fint_{B_{2R}}v^{q_0}\,\mathrm d\lambda_d\right)
 \left(\fint_{B_{2R}}v^{-q_0}\,\mathrm d\lambda_d\right)
 \le C(d,\Lambda/\lambda).
\end{equation}
\end{wniosek}
\begin{proof}
Bierzemy $g=\log v$. Z lematu~\ref{lem:pde17-log}
i nierówności Poincarégo wynika, że średnia oscylacja $g$
na każdej podkuli $B_{2R}$ jest jednostajnie ograniczona.
Uściślijmy użycie kul. Dla dowolnej
$B_\rho(y)\subset B_{3R}(x_0)$ i $\rho\le R$
kula $B_{2\rho}(y)$ mieści się w $B_{4R}(x_0)$,
więc testowanie logarytmu na podwójnej kuli daje
jednostajną kontrolę oscylacji. Jeśli $R<\rho\le3R$, kula $B_\rho(y)$ mieści się w
$B_{3R}(x_0)$. W dowodzie lematu~\ref{lem:pde17-log} wybieramy teraz
odcięcie równe $1$ na $B_{3R}(x_0)$, o zwartym nośniku
w $B_{4R}(x_0)$ i $|\nabla\eta|\le C_d/R$. Daje to
$\int_{B_{3R}(x_0)}|\nabla g|^2\,\mathrm d\lambda_d(x)
\le C R^{d-2}$.
Nierówność Poincarégo na $B_\rho(y)$ i $\rho>R$
prowadzą zatem do
$\frac{1}{\lambda_d(B_\rho(y))}\int_{B_\rho(y)}|g(x)-g_{B_\rho(y)}|\,\mathrm d\lambda_d(x)
\le C(d,\Lambda/\lambda)(R/\rho)^{(d-2)/2}
\le C(d,\Lambda/\lambda)$.
W obu przypadkach
otrzymujemy ograniczenie średnich oscylacji zależne
tylko od $d$ i $\Lambda/\lambda$. Stosujemy zatem
zlokalizowaną wersję lematu Johna--Nirenberga z kulą
mniejszą $B_{2R}(x_0)$ i większą $B_{3R}(x_0)$.
Jako $M$ przyjmujemy dodatnią stałą nie mniejszą niż $1$
i nie mniejszą niż uzyskane ograniczenie średnich oscylacji.
W lemacie~\ref{lem:pde17-jn} dobieramy $q_0=c_d/M>0$.
Zachodzą tożsamości
$v^{\pm q_0}=\exp(\pm q_0(g-g_{B_{2R}}))
\exp(\pm q_0g_{B_{2R}})$.
Mnożąc całki i używając~\eqref{eq:pde17-jn}, usuwamy
przeciwne czynniki zależne od średniej $g$.
\end{proof}

\section{Nierówność Harnacka i regularność Höldera}
\label{sec:pde17-harnack}
\begin{twierdzenie}[Nierówność Harnacka dla słabych rozwiązań]
\label{thm:pde17-harnack}
Niech $d\ge2$, $A$ spełnia~\eqref{eq:pde17-ellipticity},
$B_{4R}(x_0)\subset\Omega$, zaś $u\in H^1(B_{4R})$
będzie nieujemnym słabym rozwiązaniem
$-\operatorname{div}(A\nabla u)=0$.
Wówczas istnieje stała $C=C(d,\Lambda/\lambda)$ taka, że
\begin{equation}\label{eq:pde17-harnack}
 \operatorname*{ess\,sup}_{B_R(x_0)}u
 \le C\operatorname*{ess\,inf}_{B_R(x_0)}u.
\end{equation}
\end{twierdzenie}
\begin{proof}
Pracujemy na obszarze $B_{4R}(x_0)$; wszystkie funkcje
testowe mają zwarty nośnik w tej kuli. W szczególności
$\overline{B_{2R}(x_0)}\subset B_{4R}(x_0)$, więc
$v=u+\eps$ jest ograniczona na $B_{2R}$ i spełnia
hipotezy lematu~\ref{lem:pde17-moser-iteration}.
Do $v$ stosujemy kolejno iterację
\eqref{eq:pde17-moser-up} dla potęgi dodatniej i
\eqref{eq:pde17-moser-inv} dla potęgi odwrotnej,
z wykładnikiem $q_0$ z~\eqref{eq:pde17-bridge}:
\[
 \operatorname*{ess\,sup}_{B_R}v
 \le C\left(\fint_{B_{2R}}v^{q_0}\,\mathrm d\lambda_d\right)^{1/q_0},
 \qquad
 \left(\fint_{B_{2R}}v^{-q_0}\,\mathrm d\lambda_d\right)^{-1/q_0}
 \le C\operatorname*{ess\,inf}_{B_R}v.
\]
Przy przejściu od norm do średnich wykorzystujemy
$R^{-d/q_0}|B_{2R}|^{1/q_0}=|B_2|^{1/q_0}$;
w drugim oszacowaniu odwracamy dodatnie wielkości
i stosujemy tożsamość
$\operatorname*{ess\,sup}v^{-1}=1/\operatorname*{ess\,inf}v$.
Oznaczmy $P=\fint_{B_{2R}}v^{q_0}\,\mathrm d\lambda_d$,
$N=\fint_{B_{2R}}v^{-q_0}\,\mathrm d\lambda_d$.
Obie wielkości są dodatnie i skończone. Z mostu
~\eqref{eq:pde17-bridge} mamy $PN\le C_0$, skąd
$P^{1/q_0}\le C_0^{1/q_0}N^{-1/q_0}$.
Łącząc trzy oszacowania, otrzymujemy
$\operatorname*{ess\,sup}_{B_R}v\le
C\operatorname*{ess\,inf}_{B_R}v$ ze stałą niezależną od
$\eps$. Ponieważ dodanie stałej $\eps$ przesuwa zarówno
kres istotny górny, jak i dolny dokładnie o $\eps$,
przejście $\eps\downarrow0$ kończy dowód.
\end{proof}

\begin{twierdzenie}[Twierdzenie De Giorgiego--Nasha--Mosera o regularności lokalnej]
\label{thm:pde17-holder}
Niech $d\ge2$, $\Omega\subset\R^d$ będzie obszarem,
macierz $A$ spełnia~\eqref{eq:pde17-ellipticity}, a
$u\in H^1_{\rm loc}(\Omega)$ niech będzie słabym rozwiązaniem
$-\operatorname{div}(A\nabla u)=0$.
Istnieje $\alpha=\alpha(d,\Lambda/\lambda)\in(0,1)$
taka, że $u$ ma reprezentanta należącego do $C^{0,\alpha}_{\rm loc}(\Omega)$.
Dla $B_{4R}(x_0)\Subset\Omega$ zachodzi oszacowanie
\begin{equation}\label{eq:pde17-holder}
 [u]_{C^{0,\alpha}(\overline{B_R(x_0)})}
 \le C R^{-\alpha}
 \left(\fint_{B_{4R}(x_0)}|u|^2\,\mathrm d\lambda_d\right)^{1/2},
\end{equation}
z $C$ zależną tylko od $d$ i $\Lambda/\lambda$.
\end{twierdzenie}
\begin{proof}
Najpierw uzasadnimy redukcję oscylacji. Oznaczmy
$m=\operatorname*{ess\,inf}_{B_{4r}}u$,
$M=\operatorname*{ess\,sup}_{B_{4r}}u$, dla kuli
$B_{4r}\Subset\Omega$. Z lokalnego ograniczenia
oba wyrazy są skończone. Funkcje $u-m$ i $M-u$ są
nieujemnymi rozwiązaniami w $B_{4r}$, a po zastosowaniu
nierówności Harnacka z obszarem $B_{4r}$ otrzymujemy,
oznaczając
$m_r=\operatorname*{ess\,inf}_{B_r}u$,
$M_r=\operatorname*{ess\,sup}_{B_r}u$,
\[
 M_r-m\le C(m_r-m),\qquad M-m_r\le C(M-M_r).
\]
Przyjmijmy od początku $C>1$. Jeżeli $a=m_r-m$,
$b=M-M_r$, $c=M_r-m_r$, to
$c\le(C-1)a$ i $c\le(C-1)b$. Dodając, dostajemy
$2c\le(C-1)(a+b)=(C-1)(M-m-c)$, czyli
$(C+1)c\le(C-1)(M-m)$. Stąd
\begin{equation}\label{eq:pde17-osc-decay}
 \operatorname*{ess\,osc}_{B_r}u
 \le\vartheta\operatorname*{ess\,osc}_{B_{4r}}u,
 \qquad \vartheta=\frac{C-1}{C+1}<1.
\end{equation}
Z wyboru stałej Harnacka mamy $0<\vartheta<1$. Kładziemy
$\alpha=\min\{1/2,-\log_4\vartheta\}>0$.
Dla $0<\rho\le R$ wybieramy $k\ge0$ takie, że
$4^{-k-1}R<\rho\le4^{-k}R$. Iteracja
~\eqref{eq:pde17-osc-decay} oraz monotoniczność oscylacji
względem zawierania kul dają
\[
 \operatorname*{ess\,osc}_{B_\rho(z)}u
 \le\vartheta^k\operatorname*{ess\,osc}_{B_R(z)}u
 \le4^\alpha(\rho/R)^\alpha
       \operatorname*{ess\,osc}_{B_R(z)}u
\]
dla każdej kuli $B_R(z)\Subset\Omega$.

Ustalmy teraz $B_{4R}(x_0)\Subset\Omega$ i oznaczmy
$H=(\fint_{B_{4R}(x_0)}|u|^2\,\mathrm d\lambda_d)^{1/2}$.
Stosując~\eqref{eq:pde17-l2linfty} do $u$ i $-u$,
otrzymujemy $\|u\|_{L^\infty(B_{2R}(x_0))}\le C H$.
Niech $D$ będzie zbiorem punktów Lebesgue'a funkcji $u$
w $B_R(x_0)$; jego dopełnienie ma miarę zero.
Dla $x,y\in D$, $t=|x-y|>0$, rozważamy dwa przypadki.
Gdy $t<R/2$, oba punkty należą do $B_{2t}(x)$, a
$B_R(x)\subset B_{2R}(x_0)$. Wartości w punktach Lebesgue'a
leżą między kresami istotnymi na każdej kuli zawierającej
te punkty we wnętrzu (wynika to z granicy średnich po
małych kulach). Zatem
\[
 |u(x)-u(y)|
 \le\operatorname*{ess\,osc}_{B_{2t}(x)}u
 \le C(t/R)^\alpha H.
\]
Gdy $t\ge R/2$, korzystamy z
$|u(x)-u(y)|\le2CH\le 2^{1+\alpha}C(t/R)^\alpha H$.
Przypadek $x=y$ jest natychmiastowy.

Zbiór $D$ jest gęsty w $\overline{B_R(x_0)}$, ponieważ
każda niepusta otwarta podkula $B_R(x_0)$ ma dodatnią
miarę. Dla dowolnego punktu domknięcia wybieramy ciąg
punktów z $D$ do niego zbieżny. Powyższa nierówność
zapewnia, że ciąg wartości jest ciągiem Cauchy'ego w
$\R$ oraz że jego granica nie zależy od wyboru ciągu.
Definiujemy reprezentanta przez te granice; przechodząc
do granicy w nierówności dla par punktów z $D$,
otrzymujemy~\eqref{eq:pde17-holder} na domkniętej kuli.
Reprezentant zgadza się z $u$ prawie wszędzie.
Na przecięciach otwartych kul reprezentanci pokrywają się,
gdyż są ciągłe i równe prawie wszędzie, a zbiór pełnej
miary w otwartym przecięciu jest gęsty. Kule $B_R(x_0)$ takie, że $B_{4R}(x_0)\Subset\Omega$,
pokrywają $\Omega$. Wybieramy z nich przeliczalne
podpokrycie, korzystając z drugiego aksjomatu przeliczalności
$\R^d$. Sklejona funkcja jest równa $u$ prawie wszędzie
na całej $\Omega$, ponieważ przeliczalna suma zbiorów
miary zero ma miarę zero. Jest ciągła i lokalnie Hölderowska
z tym samym wykładnikiem; jej jednoznaczność wynika ponownie
z gęstości zbioru pełnej miary w obszarze.
\end{proof}

\begin{wniosek}[Mocna zasada maksimum przy mierzalnych współczynnikach]
Jeżeli $u$ jest ciągłym reprezentantem słabego rozwiązania
w spójnym obszarze $\Omega$ i osiąga w punkcie $\Omega$
maksimum lub minimum globalne na $\Omega$, to jest stałe
w $\Omega$.
\end{wniosek}
\begin{proof}
Niech $M=\max_\Omega u$ i $S=\{x\in\Omega:u(x)=M\}$.
Zbiór $S$ jest niepusty i względnie domknięty z ciągłości.
Dla każdego $z\in S$ wybieramy $B_{4r}(z)\Subset\Omega$.
Funkcja $M-u$ jest nieujemnym rozwiązaniem w tej kuli
i przyjmuje w $z$ wartość zero. Z ciągłości i dodatniej
miary małych kul wynika, że jej kres istotny dolny
na $B_r(z)$ jest zerowy. Nierówność Harnacka daje zerowy
kres istotny górny, a ciągłość wymusza $u=M$ na całej
$B_r(z)$. Zatem $S$ jest także otwarty. Spójność daje
$S=\Omega$; minimum rozpatrujemy, stosując wynik do $-u$.
\end{proof}

\section[Współczynniki skokowe i nieciągłość pochodnych]{Współczynniki skokowe: dlaczego nie otrzymujemy $C^1$?}
\label{sec:pde17-counterexample}
\begin{kontrprzyklad}[Brak różniczkowalności pomimo jednostajnej eliptyczności]
Dla $x\in(-1,1)$ przyjmijmy $a(x)=1$ gdy $x<0$ i
$a(x)=2$ gdy $x\ge0$. Funkcja
$u(x)=x$ dla $x<0$ oraz $u(x)=x/2$ dla $x\ge0$ należy do
$H^1(-1,1)\cap C^{0,1}([-1,1])$, ale nie do $C^1(-1,1)$.
Słaba pochodna $u'$ jest równa $1$ na $(-1,0)$ i $1/2$
na $(0,1)$: całkując przez części na obu przedziałach,
nie otrzymujemy skoku wartości $u$ w zerze, ponieważ $u(0)=0$.
Jednostronne pochodne w zerze są różne.
Ponieważ $a(x)u'(x)=1$ prawie wszędzie, mamy
$-(au')'=0$ w sensie dystrybucji. W wymiarach $d\ge2$
na kostce $\Omega=(-1,1)^d$ bierzemy
$A(x)=\operatorname{diag}(a(x_1),1,\ldots,1)$ i
$U(x)=u(x_1)$. Macierz jest mierzalna i spełnia
~\eqref{eq:pde17-ellipticity} z $\lambda=1$, $\Lambda=2$.
Z twierdzenia Fubiniego wynika
$\nabla U=(u'(x_1),0,\ldots,0)$ prawie wszędzie.
Wobec $A\nabla U=(1,0,\ldots,0)$ prawie wszędzie
dla $\varphi\in C_c^\infty(\Omega)$ mamy
$\int_\Omega A\nabla U\cdot\nabla\varphi\,\mathrm d\lambda_d=0$,
ponownie z twierdzenia Fubiniego i jednowymiarowego
twierdzenia podstawowego rachunku całkowego.
\end{kontrprzyklad}

\begin{kontrprzyklad}[Ostrość wykładnika Höldera dla macierzy mierzalnych]
\label{ex:pde17-radial}
Niech $0<\gamma<1$ i $\Omega=B_1(0)\subset\R^2$. Dla $x\ne0$
niech $e_r=x/|x|$, $e_\theta=(-x_2,x_1)/|x|$ i
\[
 A(x)=e_r\otimes e_r+\gamma^2 e_\theta\otimes e_\theta,
 \qquad u(r,\theta)=r^\gamma\cos\theta.
\]
Definiujemy $A(0)=I$. Wektory $e_r,e_\theta$ tworzą
bazę ortonormalną $\R^2$, w której macierz $A$ ma wartości
własne $1$ i $\gamma^2$. Stąd $\gamma^2I\le A\le I$,
a macierz jest mierzalna (ciągła poza zerem).
Wykażemy, że $u\in H^1(B_1)$ i jest słabym rozwiązaniem
$-\operatorname{div}(A\nabla u)=0$ na całej kuli.
Rzeczywiście, dla $r>0$ rachunek biegunowy daje
\[
 \operatorname{div}(A\nabla u)=
 \frac1r\partial_r(r\partial_ru)
 +\frac{\gamma^2}{r^2}\partial_{\theta\theta}u
 =\gamma^2r^{\gamma-2}\cos\theta
 -\gamma^2r^{\gamma-2}\cos\theta=0.
\]
Definiujemy $u(0)=0$. Dla $r>0$ mamy
\[
 \nabla u=\partial_r u\,e_r+r^{-1}\partial_\theta u\,e_\theta
 =r^{\gamma-1}(\gamma\cos\theta\,e_r-\sin\theta\,e_\theta).
\]
Z ortonormalności wynika $|\nabla u|\le r^{\gamma-1}$,
zatem
\[
 \int_{B_1}|\nabla u|^2\,\mathrm d\lambda_2
 \le 2\pi C^2\int_0^1 r^{2\gamma-1}\,\mathrm dr<\infty.
\]
Ponadto $|u|\le1$. Aby utożsamić gradient klasyczny poza
zerem z gradientem słabym, całkujemy przez części na
$B_1\setminus\overline{B_\eps}$ z testem
$\varphi\in C_c^\infty(B_1)$. Wyraz brzegowy
dla pierwszej pochodnej jest ograniczony przez
$2\pi\eps^{1+\gamma}\|\varphi\|_\infty$, więc znika.
Całki objętościowe przechodzą do granicy z całkowalności
$u$ i gradientu. To dowodzi $u\in H^1(B_1)$.
W analogicznym rachunku dla strumienia $A\nabla u$
wyraz brzegowy ma moduł co najwyżej
$2\pi C\eps^\gamma\|\varphi\|_\infty$, a całka
objętościowa z dywergencją poza zerem jest zerowa.
Przejście do granicy daje~\eqref{eq:pde17-weak}.

Sprawdźmy dokładnie regularność Höldera. Dla
$x,y\in B_1$ przyjmijmy $r=|x|\ge|y|$ i $t=|x-y|>0$.
Jeśli $t\ge r/2$, to
$|u(x)-u(y)|\le2r^\gamma\le 2^{1+\gamma}t^\gamma$.
Jeśli $t<r/2$, każdy punkt odcinka $[x,y]$
ma normę co najmniej $r-t>r/2$. Z oszacowania
gradientu i całkowania wzdłuż odcinka otrzymujemy
$|u(x)-u(y)|\le C (r/2)^{\gamma-1}t\le C t^\gamma$.
Zatem $u\in C^{0,\gamma}(B_1)$.
Nie należy do $C^{0,\beta}$ w żadnym otoczeniu zera
dla $\beta>\gamma$, gdyż dla $x=(r,0)$ mamy
$|u(x)-u(0)|/|x|^\beta=r^{\gamma-\beta}\to\infty$.
W szczególności nie istnieje jeden dodatni wykładnik
Höldera obowiązujący dla wszystkich ilorazów
eliptyczności: w tym przykładzie $\Lambda/\lambda=\gamma^{-2}$.
Nie twierdzimy, że przykład wyznacza optymalny wykładnik
dla każdej klasy macierzy w dowolnym wymiarze.
\end{kontrprzyklad}

\section{Granice teorii i dalsze kierunki}
\label{sec:pde17-final}
Dowiedliśmy nierówności Harnacka oraz regularności Höldera
dla jednorodnych, \emph{skalarnych} równań dywergencyjnych
przy jedynie mierzalnych i jednostajnie eliptycznych
współczynnikach. Kluczowy mechanizm to oszacowanie energii,
zanurzenie Sobolewa, iteracja i kontrola $\log u$.
Dla $-\operatorname{div}(A\nabla u)=f$ prawa strona wnosi
nowe człony: typowym warunkiem regularności jest
$f\in L^q_{\rm loc}$ z $q>d/2$ (dla odpowiednio dobranej
klasy rozwiązań). W twierdzeniach dla układów równań
sytuacja może być zasadniczo gorsza. Rozdział kolejny
zajmie się teorią spektralną, ale osobno zachowamy miejsce
na pełne warianty niejednorodne, paraboliczne i kontrprzykłady.

% -----------------------------------------------------------------------
% Rozdział 18 — do włączenia PRZED \appendix w pliku wykładu
% -----------------------------------------------------------------------
